% Review-language results (M.N.). Planned content:
%  - how reviewers use "authentic" (concordance; share affirming / negating),
%    validated on a hand-labelled sample;
%  - a review-based score per recipe (authentic mentions + sentiment),
%    controlling for the number of reviews;
%  - Fig. 3: the three ingredient scores vs. the review score, per cuisine.

\subsection*{What reviewers call authentic}
TODO (NLP): how often and how reviewers use ``authentic''. Recipes with at least one review mentioning it have a median of 9 reviews against 2 for the rest, so the comparison must account for the number of reviews.

\subsection*{Do ingredients predict perceived authenticity?}
TODO: compare the three ingredient scores with the review score (Fig.~\ref{fig:comparison}). Which of \emph{what}, \emph{which pairs} or \emph{how} predicts it best, and in which cuisines?

\begin{figure}[tbhp]
\centering
\fbox{\parbox[c][5cm][c]{0.9\linewidth}{\centering\sffamily\small
TODO Fig.~3: ingredient scores vs. review-based authenticity,\\
one panel or colour per cuisine, against a null.}}
\caption{TODO. The answer to the research question.}
\label{fig:comparison}
\end{figure}
